%%%%%%%%%%%%%%%%%%%%%%%%%%%%%%%%%%%%%%%%%%%
% Página de resumen del proyecto en inglés%
%%%%%%%%%%%%%%%%%%%%%%%%%%%%%%%%%%%%%%%%%%%
% !TEX root = A0.MiTFG.tex

\pagestyle{fancy}
\addstarredchapter{Abstract}

\begin{center}
	\scshape
	E.T.S. de Ingeniería de Telecomunicación, Universidad de Málaga
\end{center}

\bigskip

\begin{center}
	\Large \scshape
	\textbf{\tfgtitlenameENG}
\end{center}

\bigskip \bigskip \bigskip

\begin{minipage}{\textwidth}

\textbf{Author:} \tfgauthorname

\medskip

\textbf{Supervisor:} \tfgtutorname

%\textbf{Co-supervisor:} <Nombre del cotutor>\ (elimina esta línea si no hay cotutor)

\medskip

\textbf{Department:} Electronic Technology

\medskip

\textbf{Degree:} Degree in Electronic Systems Engineering

\medskip

\textbf{Keywords:} RISC-V, FPGA, CPU architecture, open source hardware, digital design, simulation, HDL, Verilog, embedded processor, system-on-chip, cache, pipeline, functional validation, hardware development.

\bigskip \bigskip


\end{minipage}

\begin{center}
	\textbf{Abstract}
\end{center}

The present project aims to design a processor based on the RISC-V architecture, specifically implementing the RV32I instruction set. To ensure its correct operation, a verification method based on signature generation and comparison is employed. Finally, the processor is implemented on an FPGA platform in order to experimentally validate its performance and verify compliance with the architectural specification.

\blankpage
